\section{Local uniqueness tests for the complete arithmetic operator (v1.55)}
\label{sec:v155-local-uniqueness}

The affine-potential criterion of Proposition~\ref{prop:ns51-affine-potential}
asks for the absence of a nullvector satisfying the equation on the
\emph{whole} window. Here we test three possible mechanisms for that
uniqueness. The complete arithmetic operator fails a local-patch
unique-continuation property, and its semigroup does not preserve the
ordinary cone of nonnegative functions. Neither statement gives a
nullvector on the whole window or a negative diagonal Weil direction.
An exact inward-dilation identity isolates the arithmetic overlap
terms still required at a hypothetical first zero. API, the uniform
floor, G2 and RH remain open. No numerical window or tail metric is added.

% NS-55: exact and scoped. Local-patch witnesses are not full nullvectors.
\subsection*{Unique continuation: the arithmetic shifts change the theorem}

Use $I_a=(-a,a)$, its complete operator $A_a$, and
\[
 \ell=\log2,\quad w_2=\frac{\log2}{\sqrt2},\quad
 \rho(t)=\frac{e^{-t/2}}{1-e^{-2t}},\quad
 R(t)=2\cosh(t/2)-\rho(t)\quad(t>0).
\]
If $f$ vanishes on an open $J\subset I_a$ separated from its support,
the exact operator there is
\begin{equation}\label{eq:ns55-off-support-operator}
 (A_af)(x)=\int_{I_a}R(|x-y|)f(y)\,dy
 -\sum_{1<n<e^{2a}}\frac{\Lambda(n)}{\sqrt n}
       \{E_af(x-\log n)+E_af(x+\log n)\}.
\end{equation}
This follows from Proposition~\ref{prop:v118-log-dirichlet}: local
diagonal terms vanish, and the logarithmic kernel and its correction
combine into $-\rho$. Equivalently it is the off-origin distribution
identity for $-\mathfrak g''$. In particular the pole term is retained.

\begin{proposition}[Failure of local open-set unique continuation]
\label{prop:ns55-local-ucp-counterexample}
For every $a>\tfrac12\log2$, there are a nonempty $U\Subset I_a$
and a nonzero real $f\in\mathcal D(A_a)$ with
\[
 f|_U=0,\qquad (A_af)|_U=0.
\]
The second equality is pointwise and distributional on $U$ only.
No complete nullvector or negative Weil direction is asserted.
\end{proposition}
\begin{proof}
Set
\[
 \delta_0=\min\{\ell/8,\tfrac14\log(3/2)\},\qquad
 M_0=2\cosh((\ell+2\delta_0)/2)+\rho(\ell-2\delta_0).
\]
Choose $0<\epsilon<\min\{a-\ell/2,\delta_0,w_2/(4M_0)\}$ and put
\[
 U=(\ell/2-\epsilon,\ell/2+\epsilon),\quad V=U-\ell,\quad
 W=(-\epsilon,\epsilon).
\]
Their closures are disjoint and inside $I_a$. Choose any nonzero real
$h\in C_c^\infty(W)$ and define on $L^2(U)$
\[
 (Tv)(x)=\int_U R(x-t+\ell)v(t)\,dt,\qquad
 H(x)=\int_W R(x-y)h(y)\,dy .
\]
Since $|R(t)|\le M_0$ for $|t-\ell|\le2\delta_0$, Schur's estimate
gives $\|T\|_{2\to2}\le2\epsilon M_0<w_2/2$. Define
$v=(w_2I-T)^{-1}H$ by its convergent Neumann series. The equation
$w_2v=H+Tv$ makes $v$ smooth on a neighborhood of $\overline U$:
the smooth kernels and their derivatives are bounded on that
neighborhood times the fixed integration intervals.

Set $f(y)=v(y+\ell)$ on $V$, $f=h$ on $W$, and zero elsewhere.
This is nonzero. On $U$, positive prime shifts lie beyond the support.
For $n\ge3$, $2\epsilon<\log(3/2)$ puts $x-\log n$ left of $V$.
The only nonzero translated value can therefore be $n=2$, equal to
$v(x)$. Equation~\eqref{eq:ns55-off-support-operator} gives
$A_af=-w_2v+Tv+H=0$ on $U$.

The zero extension is compactly supported and piecewise smooth with
finitely many jumps, hence belongs to $H^s(\mathbb R)$ for
$0<s<1/2$. Its full-line logarithmic multiplier lies in $L^2$:
$\log^2|\xi|\ll_s1+|\xi|^{2s}$ at high frequency, and its transform
is bounded near zero, where $\log^2|\xi|$ is integrable. The form
representation therefore puts $f$ in
$\mathcal D(\tfrac12L_\Delta^{\rm Dir})$. Bounded perturbation
gives $f\in\mathcal D(A_a)$. This is the complete operator, not a head.
\end{proof}

\begin{proposition}[Parity and finitely many linear constraints]
\label{prop:ns55-parity-ucp-counterexample}
For every $a>\log2$ and $\sigma\in\{1,-1\}$, there are a nonempty
$U\Subset(0,a)$ and an infinite-dimensional real linear subspace
$\mathcal L_\sigma\subset\mathcal D(A_a)$ such that
\[
 f(-x)=\sigma f(x),\qquad f=A_af=0\text{ on }U\cup(-U)
 \quad(f\in\mathcal L_\sigma).
\]
One may impose any prescribed finite number of real linear conditions
on $\mathcal L_\sigma$ and still obtain nonzero witnesses. In particular,
source orthogonality does not restore this local unique-continuation
statement. This assertion remains about a local patch, not the full kernel.
\end{proposition}
\begin{proof}
Put $\mathcal P_a=\{\log n:1<n<e^{2a},\ \Lambda(n)>0\}$.
Choose $v_0\in(0,a-\ell)$ avoiding the finitely many equations
$2v_0+\ell\in\mathcal P_a$, and let $u_0=v_0+\ell$.
Choose $w_0\in(0,v_0/2)$ with
$u_0-w_0,u_0+w_0\notin\mathcal P_a$.

Choose an initial radius $\epsilon_0>0$ such that the intervals
$U_{\epsilon_0},V_{\epsilon_0},W_{\epsilon_0}$ about $u_0,v_0,w_0$
and their reflections have disjoint closures inside $I_a$.
Choose it smaller if necessary so differences from $U_{\epsilon_0}$
to $V_{\epsilon_0},-V_{\epsilon_0},W_{\epsilon_0},-W_{\epsilon_0}$
meet $\mathcal P_a$ only in the intended $U-V$ difference $\ell$.
There are only finitely many forbidden central differences, all
separated from $\mathcal P_a$, so this is possible. Set
\[
 M_\sigma=\sup_{x,t\in\overline{U_{\epsilon_0}}}
  |R(x-t+\ell)+\sigma R(x+t-\ell)|<\infty .
\]
Choose $0<\epsilon\le\epsilon_0$ with
$2\epsilon M_\sigma<w_2/2$, and write $U,V,W$ for the smaller intervals.
For arbitrary $h\in C_c^\infty(W)$, define
\begin{align*}
 (T_\sigma v)(x)&=\int_U
       \{R(x-t+\ell)+\sigma R(x+t-\ell)\}v(t)\,dt,\\
 H_\sigma(x)&=\int_W
       \{R(x-y)+\sigma R(x+y)\}h(y)\,dy .
\end{align*}
All arguments are positive. Solve
$(w_2I-T_\sigma)v=H_\sigma$ by the same norm bound.
Set $f(y)=v(y+\ell)$ on $V$, $f=h$ on $W$, reflect with sign $\sigma$,
and put it zero elsewhere. The same smoothness and domain proof applies.
The sole nonzero prime-shift term on $U$ is $-w_2v$, and the integrals
give $T_\sigma v+H_\sigma$, hence $A_af=0$ there. Parity gives the
reflected equality. The linear map $h\mapsto f$ is injective since
$f|_W=h$, so its image is infinite dimensional. Its intersection with
the kernels of finitely many linear functionals is still infinite
dimensional. A complex orthogonality condition on real witnesses is
at most two real linear conditions.
\end{proof}

\paragraph{The logarithmic theorem and its missing hypotheses.}
Chen--Hauer--Weth~\cite[Theorem 1.7]{ChenHauerWeth2023} prove that
$u=L_\Delta u=0$ on the \emph{same} nonempty open set implies $u=0$
on $\mathbb R$, for $u\in L^1((1+|x|)^{-1}dx)$.
See also Harrach--Lin--Weth~\cite[Proposition 4.1]{HarrachLinWeth2025}.
Compactly supported $L^2$ functions satisfy the integrability hypothesis.
But on an interior open set the arithmetic equation gives
$L_\Delta E_af=-2\mathcal K_a f$; $f=0$ does not annihilate this
nonlocal remainder. The preceding counterexamples show the difference
for the actual operator. Replacing it by a local multiplication
potential would change the problem.

Outside the window, only $E_af=0$ is prescribed; neither
$L_\Delta E_af=0$ nor the full arithmetic equation is prescribed.
For example, at $x>a$ one has
\[
 (L_\Delta E_af)(x)=-\int_{-a}^a\frac{f(y)}{x-y}\,dy,
\]
which generally is nonzero. Endpoint trace zero supplies no open
interior vanishing set or missing exterior equation.

\begin{proposition}[Uniqueness with a prime-shift-free exterior-of-support cell]
\label{prop:ns55-prime-free-cell}
Let $f\in V_a$ have essential support $S\subset[-a,a]$, with
$b=\sup S<a$. Suppose a nonempty open interval $J\Subset(b,a)$ obeys
\[
 J\cap\bigcup_{s\in\mathcal P_a}(S+s)=\varnothing,\qquad
 A_af=0\text{ on }J\text{ in distributions}.
\]
Then $f=0$. The reflected left-hand statement also holds.
Here $A_af$ is understood weakly if only $f\in V_a$ is initially given.
\end{proposition}
\begin{proof}
Assume $S$ nonempty. Positive shifts lie beyond $S$; the displayed
separation removes negative shifts. Thus
\[
 \mathcal R_f(x)=\int_{I_a}R(x-y)f(y)\,dy=0\quad(x\in J).
\]
This is real analytic for all $x>b$, so vanishes on that half-line.
Put
$M_-=\int e^{-y/2}f(y)\,dy$ and
$M_j=\int e^{(2j+1/2)y}f(y)\,dy$.
The locally uniformly convergent identity
$\rho(t)=\sum_{j\ge0}e^{-(2j+1/2)t}$ for $t>0$ gives
\[
 0=e^{-x/2}\mathcal R_f(x)
   =M_- -\sum_{j=1}^\infty M_j e^{-(2j+1)x}\quad(x>b).
\]
The $j=0$ term cancels exactly one pole contribution. Since
$|M_j|\le\|f\|_1e^{(2j+1/2)b}$, the last series is a power
series in $z=e^{-x}$ with radius at least $e^{-b}$.
Taking $x\to\infty$ and then its coefficients gives $M_-=0$
and $M_j=0$ for every $j\ge1$.

Push the finite complex measure $e^{y/2}f(y)\,dy$ forward by
$t=e^{2y}$. It is supported in $[e^{-2a},e^{2b}]$, away from zero,
and all its positive integer moments vanish. Polynomials without a
constant term are dense among continuous functions on this interval:
approximate the desired function divided by $t$ and multiply by $t$.
Thus the measure is zero, proving $f=0$.
Only the separated integral component was analytically continued,
not the whole arithmetic equation with its shifts.
\end{proof}

\paragraph{Scope of the remaining gap.}
The last proposition requires actual support separation, not smallness
or zero endpoint trace. A separate first-crossing support argument forces
a first-crossing nullvector to reach both window endpoints; hence it
cannot supply the last proposition's exterior-of-support cell inside
the window. The local counterexamples are not first-crossing nullvectors.
They disprove the direct transplantation of the logarithmic open-set
unique-continuation theorem to the full shifted operator, while the
prime-free-cell theorem is a restricted arithmetic uniqueness result.
Complete affine-potential injectivity, G2 and RH remain open.

% NS-55 first-crossing subtask. Exact identities and scoped countermodel.
% No H^1 assumption, shape derivative, numerical window, or RH assumption.
\subsection*{First-crossing geometry on the complete form domain}

Let $V_a=\mathcal D(q_a)$ be the complete form domain from NS-51,
with unitary Fourier transform of the zero extension understood.
Let $\mu(a)=\inf\sigma(A_a)$. A \emph{first-zero window} means
\begin{equation}\label{eq:ns55fc-first}
 \mu(a_*)=0,\qquad \mu(b)>0\quad(0<b<a_*).
\end{equation}
This is a conditional setup, not an assertion that such a window exists.
If a zero or negative lower edge occurs, continuity and small-window
positivity yield such a first window. The required continuity is
Suzuki's Theorem 1.3, with its complete-domain scaling argument in
Section 4 of version 2
(\url{https://arxiv.org/html/2606.09096v2#S4}).

\begin{proposition}[Support saturation and exact inward variation]
\label{prop:ns55fc-geometry}
Assume \eqref{eq:ns55fc-first}, and let $0\ne f\in\ker A_{a_*}$.
Then the essential support of its zero extension has diameter $2a_*$;
its extreme support points are $-a_*$ and $a_*$. For $0<r<1$ set
\[
 (U_rf)(x)=r^{-1/2}f(x/r),
\]
where $f$ has already been extended by zero. Then $U_rf\in V_{ra_*}$,
$\|U_rf\|_2=\|f\|_2$, and
\begin{equation}\label{eq:ns55fc-null-variation}
 q_{ra_*}[U_rf]
 =q_{a_*}[U_rf-f]
 \ge\mu(ra_*)\|f\|_2^2>0.
\end{equation}
All identities are valid on the complete form domain. Inward dilation
preserves each parity sector.
\end{proposition}
\begin{proof}
The global compact-support form is invariant under simultaneous
translation, and its restrictions are support consistent. If the
essential support diameter were smaller than $2a_*$, a translation
would place $f$ in $V_b$ for some $b<a_*$. Its unchanged form value
zero contradicts $q_b\ge\mu(b)\|\cdot\|_2^2$.

The Fourier identity $\widehat{U_rf}(\xi)=\sqrt r\,\widehat f(r\xi)$
and comparison of $\log(2+|\xi|/r)$ with $\log(2+|\xi|)$ prove
membership in $V_{ra_*}$. Since $f$ is an operator nullvector,
$q_{a_*}(f,v)=0$ for every $v\in V_{a_*}$ by the representation
theorem for closed forms. Expand $q_{a_*}[U_rf-f]$; both cross terms
and $q_{a_*}[f]$ vanish. Support consistency and the variational
principle on $I_{ra_*}$ give \eqref{eq:ns55fc-null-variation}.
No derivative of $r\mapsto U_rf$ is taken.
\end{proof}

The endpoint conclusion means that an argument requiring an empty
exterior support strip \emph{inside} $I_{a_*}$ cannot be applied to
a first-zero nullvector without a new arithmetic step. It does not
mean that $f$ has a nonzero trace: the archived continuous zero
extension and logarithmic endpoint decay are fully compatible with
support saturation.

\begin{proposition}[Exact dilation defect with both poles and all primes]
\label{prop:ns55fc-dilation}
For $f\in V_a$, let
\[
 R_f(t)=\int_{\mathbb R}f(x+t)\overline{f(x)}\,dx,
 \quad P_f(r)=\int_{-a}^a f(x)\cosh(rx/2)\,dx,
 \quad S_f(r)=\int_{-a}^a f(x)\sinh(rx/2)\,dx,
\]
and set $\alpha(\xi)=\operatorname{Re}\psi(1/4+i\xi/2)-\log\pi$.
For $0<r<1$ define
\begin{align}
 \mathcal A_r[f]&=\int_{\mathbb R}
    [\alpha(\xi/r)-\alpha(\xi)]|\widehat f(\xi)|^2\,d\xi,
       \label{eq:ns55fc-arch}\\
 \mathcal P_r[f]&=2r\bigl(|P_f(r)|^2-|S_f(r)|^2\bigr)
                  -2\bigl(|P_f(1)|^2-|S_f(1)|^2\bigr),\notag\\
 \mathcal R_r[f]&=2\sum_{1<n<e^{2a}}\frac{\Lambda(n)}{\sqrt n}
          \operatorname{Re}\{R_f(\log n/r)-R_f(\log n)\}.
          \notag
\end{align}
Then
\begin{equation}\label{eq:ns55fc-defect}
 q_{ra}[U_rf]-q_a[f]
       =\mathcal A_r[f]+\mathcal P_r[f]-\mathcal R_r[f].
\end{equation}
For nonzero $f$,
\begin{equation}\label{eq:ns55fc-arch-bound}
 0<\mathcal A_r[f]\le5\log(1/r)\|f\|_2^2.
\end{equation}
In the even sector $S_f(r)=0$; in the odd sector $P_f(r)=0$.
Neither pole sign may be dropped in the full-space identity.
\end{proposition}
\begin{proof}
The exact complete formula, obtained from
\eqref{eq:v129-jump-decomposition} and
Proposition~\ref{prop:v135-both-parities}, is
\[
 q_a[f]=\int\alpha(\xi)|\widehat f(\xi)|^2\,d\xi
  +2|P_f(1)|^2-2|S_f(1)|^2
  -2\sum_{1<n<e^{2a}}\frac{\Lambda(n)}{\sqrt n}
                        \operatorname{Re}R_f(\log n).
\]
This extends from the smooth core by the form norm: at fixed $a$
the pole and finite prime operators are bounded, and the logarithmic
multiplier defines the principal closed form. Change variables in
the Fourier integral. Directly,
$R_{U_rf}(t)=R_f(t/r)$ and
$P_{U_rf}(1)=\sqrt r P_f(r)$, with the same formula for $S$.
The correlation $R_f$ is continuous, and $R_f(t)=0$ for
$|t|\ge2a$. Thus the dilated prime sum can use the fixed upper
cutoff $e^{2a}$: every added term past $e^{2ra}$ is exactly zero,
including equality. This proves \eqref{eq:ns55fc-defect}, retaining
every changed support overlap rather than differentiating the prime
cutoff or discarding boundary terms.

For $s_k=k+1/4$ and $t=|\xi|/2$, the convergent digamma series gives
\[
 \xi\alpha'(\xi)
 =\sum_{k\ge0}\frac{2t^2s_k}{(s_k^2+t^2)^2}>0\quad(\xi\ne0).
\]
For fixed $t>0$, $F_t(s)=2t^2s/(s^2+t^2)^2$ is nonnegative and
unimodal. Comparing its unit-spaced samples on either side of its
maximum with their adjacent integrals gives
\[
 \sum_{k\ge0}F_t(k+1/4)
 \le\int_{1/4}^\infty F_t(s)\,ds
               +2\sup_{s\ge1/4}F_t(s)\le1+4=5;
\]
the last estimate uses $F_t(s)\le1/(2s)$. Integrate this logarithmic
derivative from $|\xi|$ to $|\xi|/r$. Strict increase away from the
origin and the fact that a nonzero $L^2$ transform cannot be supported
at one point prove \eqref{eq:ns55fc-arch-bound}.
\end{proof}

At a first-zero nullvector, the exact necessary inequality is therefore
\begin{equation}\label{eq:ns55fc-necessary}
 \mathcal R_r[f]
 \le\mathcal A_r[f]+\mathcal P_r[f]
             -\mu(ra_*)\|f\|_2^2
 <\mathcal A_r[f]+\mathcal P_r[f],\qquad 0<r<1.
\end{equation}
There is no contradiction: inward scaling has positive energy.
An independent arithmetic estimate forcing the reverse non-strict
inequality for even one $r<1$ at every proposed first-zero nullvector
would exclude such a nullvector. No such estimate is derived here.
Merely naming that reverse inequality restates a sufficient exclusion
test; it does not make an arithmetic lemma out of the variational identity.

The form domain does not in general make $R_f$ differentiable at the
prime shifts. Its Fourier representation controls a logarithmic moment,
not $\int|\xi||\widehat f(\xi)|^2d\xi$. Consequently dividing
\eqref{eq:ns55fc-defect} by $1-r$ and passing to a differentiated
virial identity would require an additional regularity or limit
argument. The finite identity requires neither. Strong continuity of
dilations in the logarithmic form topology, and hence continuity of
the defect, follows by smooth approximation and locally uniform
dilation bounds in that topology; it supplies no differentiability rate.


% NS-55 prime-shift lane: exact identities and scoped operator obstructions.
% Uses v118 physical realization and v154 affine-potential notation.
\subsection*{Prime-shift primitives and the limits of transfer arguments}

\paragraph{Scope.}
Let $I_a=(-a,a)$, $L=2a$, and use the complete operator $A_a$,
form domain $V_a$ and screw function $\mathfrak g$ of
Proposition~\ref{prop:ns51-affine-potential}. Write
$\ell_n=\log n$ and $w_n=\Lambda(n)/\sqrt n$. Terms with
$w_n=0$ are omitted. The following are exact identities and proved
obstructions to particular inferences. They prove neither an arithmetic
nullvector nor affine-potential injectivity (API). Both pole signs are
retained. No new window or tail metric is computed; G2 and RH remain open.

\begin{proposition}[The exact two-sided primitive equation]
\label{prop:ns55-prime-primitives}
Remove the prime ramps from the screw function, writing on $|t|<L$
\[
 \mathfrak g(t)=\mathfrak g_0(t)+
       \sum_{1<n<e^L}w_n(|t|-\ell_n)_+.
\]
Thus $\mathfrak g_0$ includes the archimedean and both pole terms.
For $f\in L^2(I_a)$ define the clamped primitive and total integral
\[
 H_f(x)=\int_{-\infty}^x E_af(y)\,dy,
 \qquad M_f=\int_{I_a}f(y)\,dy.
\]
Then, at every $x\in I_a$,
\begin{equation}\label{eq:ns55-prime-primitive}
 U_f(x)=(\mathfrak g_0'*E_af)(x)+
 \sum_{1<n<e^L}w_n
       \{H_f(x-\ell_n)+H_f(x+\ell_n)-M_f\}.
\end{equation}
In particular API's equation $U_f=b$ is exactly this two-sided
shifted-primitive equation, subject to $f\in V_a$ and the parity constants
of Proposition~\ref{prop:ns51-affine-potential}. Differentiation is valid
in distributions and gives
\begin{equation}\label{eq:ns55-prime-differentiated}
 (\mathfrak g_0''*E_af)|_{I_a}
 +\sum_{1<n<e^L}w_n\{E_af(\cdot-\ell_n)+E_af(\cdot+\ell_n)\}|_{I_a}=0.
\end{equation}
No $H^1$ regularity of $f$ is required or inferred.
\end{proposition}
\begin{proof}
The derivative of $(|t|-\ell)_+$ is
$\mathbf1_{t>\ell}-\mathbf1_{t<-\ell}$ almost everywhere. Its
convolution with $E_af$ is
$H_f(x-\ell_n)-\{M_f-H_f(x+\ell_n)\}$.
The local $L^2$ property of $\mathfrak g_0'$ and compact support of
$E_af$ make its convolution continuous, as in the affine-potential
criterion. The finite sum is continuous too. Finally
$H_f'=E_af$ in distributions. Membership of this antiderivative in
$H^1_{\rm loc}$ does not imply $f\in H^1$.
\end{proof}

\paragraph{Why this is not a propagation rule.}
The prime part couples both $x-\ell_n$ and $x+\ell_n$.
More importantly, $\mathfrak g_0''*E_af$ remains a nonlocal term.
Knowing $f=0$ on a subinterval does not make that term zero there;
values of $f$ throughout the interval still enter. The logarithmic
singularity also lies on the convolution diagonal, so analyticity
of the kernel away from its singularities does not imply analyticity
of the potential across the support. No identity theorem can be
applied to that unsupported analytic continuation. Additional
prime-free exterior hypotheses can change this conclusion; none
is imposed in the present statement.

\begin{proposition}[The full semigroup is not positivity preserving]
\label{prop:ns55-nonpositive-semigroup}
Let $r>1$ be the unique real solution of $r^3-r-1=0$, and suppose
$L>\log r$. There are real nonnegative
$f,g\in C_c^\infty(I_a)$ with disjoint supports such that
\begin{equation}\label{eq:ns55-positive-cross}
 q_a(f,g)>0.
\end{equation}
Consequently the complete semigroup $e^{-sA_a}$ is not positivity
preserving on the full real $L^2(I_a)$ space with its ordinary
nonnegative cone: for these tests and all sufficiently small $s>0$,
\[
 \langle e^{-sA_a}f,g\rangle<0.
\]
This asserts no negative diagonal Weil energy and makes no separate
claim about order preservation after a parity restriction.
\end{proposition}
\begin{proof}
For disjoint supports the diagonal terms in the operator contribute
nothing. Away from prime-power separations, the off-diagonal integral
kernel of Proposition~\ref{prop:v118-log-dirichlet} is
\[
 J(t)=2\cosh(t/2)-\rho(t),\qquad
 \rho(t)=\frac{e^{t/2}}{e^t-e^{-t}},\quad t>0.
\]
The logarithmic Laplacian contributes $-1/(2t)$ and its bounded
correction contributes $-(\rho(t)-1/(2t))$, giving $-\rho(t)$;
the pole kernel gives the displayed $2\cosh(t/2)$.
For $u=e^t>1$,
\[
 J(t)=\frac{u^3-u-1}{\sqrt u\,(u^2-1)}.
\]
The numerator is strictly increasing on $(1,\infty)$, changes sign
once, and is positive precisely when $t>\log r$.
Choose a separation $t_0\in(\log r,L)$ outside the finite set
$\{\log n:1<n<e^L,\ \Lambda(n)\ne0\}$. Two sufficiently small
disjoint intervals inside $I_a$, at separation $t_0$, have all
cross separations in the positive region of $J$ and avoid every
prime-power separation. Choose nonzero nonnegative smooth bumps
$f,g$ supported in these intervals. All translated prime cross
terms vanish and the remaining double integral is strictly positive.
The tests belong to $\mathcal D(A_a)$, and $A_a$ is self-adjoint
and bounded below at this fixed window. Strong differentiability
of its semigroup on the operator domain gives
\[
 \langle e^{-sA_a}f,g\rangle
 =\langle f,g\rangle-s\langle A_af,g\rangle+o(s)
 =-s q_a(f,g)+o(s),
\]
which is negative for small positive $s$.
\end{proof}

\paragraph{Meaning of the order obstruction.}
A standard positive-kernel or maximum-principle transfer cannot simply
be attributed to the complete arithmetic operator. This does not
exclude a more specialized uniqueness argument, positivity of a
particular eigenfunction, a different cone, or API itself.
Positivity of the quadratic form is distinct from preservation of
pointwise positivity by its semigroup. Moreover conjugation by a
strictly positive multiplier $h$ with $h,h^{-1}\in L^\infty(I_a)$
cannot repair this: multiplication by $h$ and by $h^{-1}$ both
preserve the nonnegative cone, so positivity of the conjugated
semigroup would imply positivity of the original one. This concerns
the full signed operator, not a separate positive edge form.


\paragraph{What remains after these tests.}
The whole-window API input has not been weakened by a proved arithmetic
estimate. The local counterexamples rule out one proposed uniqueness
principle for the actual operator, while the nonpositive semigroup rules
out using its ordinary cone as a maximum-principle mechanism. The
prime-free-cell result needs support separation that a first-zero
nullvector lacks. The finite dilation identity preserves the precise
signed prime overlaps and pole changes; their required estimate remains
unproved. This is a failure of these proposed mechanisms, not a
counterexample to whole-window injectivity or Weil positivity.

The full NS-55 evidence also records the norm discontinuity of an
entering compressed prime shift, its continuous screw-ramp smoothing,
and a valid fixed-window compact Fredholm reduction. The latter leaves
the physical coupling point potentially exceptional; discreteness is
not exclusion. These supporting arguments are archived in
\path{evidence/ns55_kernel/}; no global analyticity, additional endpoint
regularity, new window, tail metric, G2 or RH result is asserted.
